% from ega4-17-fr.tex
\oldpage[IV-4]{109}

$k$-algèbre, il contient un sous-corps $K'_j$ isomorphe canoniquement au corps résiduel de $A_j$ (0, 19.6.2). On prend alors $X_i=X$, $X'_i$ somme des $\Spec(K'_j)$ et il est clair que l'on répond ainsi à la question puisque pour tout $j$ on a deux homomorphismes $K'_j\to A_j\to k(A_j)$ dont le composé est un isomorphisme.

% (continued in ega4-18-fr.tex, before its first page marker)
\section[Compléments sur les morphismes étales. Anneaux locaux henséliens et anneaux strictement locaux]{Compléments sur les morphismes étales\\Anneaux locaux henséliens et anneaux strictement locaux}
\label{section:IV.18}

Dans le présent paragraphe, nous étudions diverses propriétés spéciales aux morphismes étales. De plus, la notion de morphisme étale permet de développer de façon fort naturelle la théorie de Nagata des anneaux henséliens, ainsi que celle des anneaux strictement locaux. Ces anneaux jouent un rôle important dans de nombreux développements récents, en s'introduisant chaque fois que l'on a besoin d'un procédé de « localisation » plus fin que celui fourni par la topologie de Zariski (cf. par exemple [43], en attendant la parution du chapitre de notre Traité consacré à l'étude de la « topologie étale »).

\subsection{Une équivalence remarquable de catégories.}
\EGAReaderTarget{ega4.c04.subsec.18-1}{subsection:IV.18.1}

\begin{proposition}[18.1.1]
\label{IV.18.1.1}
Soient $S$ un préschéma, $S_0$ un sous-préschéma fermé de $S$, $X_0$ un $S_0$-préschéma lisse (resp. étale) sur $S_0$, $x_0$ un point de $X_0$. Il existe alors un voisinage ouvert $U_0$ de $x_0$ dans $X_0$, un $S$-préschéma $U$ lisse (resp. étale) sur $S$ et un $S_0$-isomorphisme
\[
  U\times_S S_0\xrightarrow{\sim}U_0.
\]
\end{proposition}

Notons que si $X_0$ est étale sur $S_0$ au point $x_0$, il est a fortiori non ramifié sur $S_0$ en ce point; si l'on a construit un $S$-préschéma $U$ lisse sur $S$ tel que $U\times_S S_0$ soit isomorphe à $U_0$, comme les fibres des morphismes $U_0\to S_0$ et $U\to S$ contenant $x_0$ sont alors isomorphes, il en résultera que $U$ est non ramifié sur $S$ au point $x_0$ (17.4.1, $d$)), donc aussi dans un voisinage de $x_0$; en remplaçant $U$ par ce voisinage, on en conclut que $U$ sera étale sur $S$. Il suffit donc de prouver la proposition lorsqu'on suppose seulement $X_0$ lisse sur $S_0$.

La question étant locale sur $S$ et sur $X_0$, on peut supposer que $S=\Spec(A)$ et $X_0=\Spec(C_0)$ sont affines, de sorte que $S_0=\Spec(A_0)$, où $A_0$ est un anneau quotient de $A$, $C_0=B_0/\mathfrak J_0$, où $B_0=A_0[T_1,\ldots,T_n]$ et $\mathfrak J_0$ est un idéal de type fini de $B_0$; enfin, $C_0$ est une $A_0$-algèbre formellement lisse pour les topologies discrètes. Soit $\mathfrak p_0$ l'idéal $\mathfrak j_{x_0}$ dans $C_0$; on a $\mathfrak p_0=\mathfrak q_0/\mathfrak J_0$, où $\mathfrak q_0$ est un idéal premier de $B_0$. Le critère jacobien (0, 22.6.4) joint à (0, 19.1.12) montre qu'il existe dans $\mathfrak J_0$ une famille de $r$ polynômes $u_i$ $(1\leq i\leq r)$ et $r$ indices $j_h$ $(1\leq h\leq r)$ tels que les images des $u_i$ dans $(\mathfrak J_0)_{\mathfrak q_0}/(\mathfrak J_0^2)_{\mathfrak q_0}$ engendrent ce $(B_0)_{\mathfrak q_0}$-module et que l'on ait
\begin{equation}
  \det\left(\frac{\partial u_i}{\partial T_{j_h}}\right)\notin\mathfrak q_0.
  \tag{18.1.1.1}\label{IV.18.1.1.1}
\end{equation}

Comme $(B_0)_{\mathfrak q_0}$ est un anneau local, il résulte du lemme de Nakayama que l'on peut supposer que les images des $u_i$ dans $(\mathfrak J_0)_{\mathfrak q_0}$ engendrent ce $(B_0)_{\mathfrak q_0}$-module, puis, en remplaçant au besoin $X_0$ par un voisinage ouvert affine de $x_0$, que les $u_i$ engendrent $\mathfrak J_0$ ($0_{\mathrm I}$, 5.2.2). Posons alors $B=A[T_1,\ldots,T_n]$; $B_0$ est donc un anneau quotient

